\setlength{\parskip}{1em}
Credit card fraud detection is a problem of severe class imbalance. In the dataset used in
this project, 473 of 283{,}726 transactions are fraudulent, a rate of 0.167 percent. Plain
accuracy is therefore useless as a measure of quality: a model that labels every
transaction as legitimate is right 99.83 percent of the time and catches no fraud at all.
This project builds an end-to-end system that trains and compares eight classifier
pipelines. Four algorithms (Logistic Regression, Random Forest, XGBoost, and LightGBM) are
each trained under two strategies for handling imbalance (class weighting and SMOTE
oversampling). Models are ranked by the area under the precision-recall curve (PR-AUC),
which is the right metric to use when the positive class is rare. To prevent leakage,
1{,}081 exact duplicate rows are removed before the data is split, and any resampling is
applied strictly inside the training fold of each pipeline, so it never touches the
validation or test data.

The best model on a single split was LightGBM with class weighting. It reached a PR-AUC of
0.828 on a held-out test set of 56{,}746 transactions, at a decision threshold of 0.636
tuned on validation data. At that operating point it catches 70 of 95 frauds while raising
only 2 false alarms.

The main contribution of the project is what happens when that headline number is tested
statistically. A 1{,}000-resample bootstrap places the champion's PR-AUC in a 95 percent
interval of $[0.751, 0.899]$, which overlaps almost entirely with the intervals of the next
four models. Pairwise McNemar tests find that none of the 28 model pairs differ
significantly after Holm correction for multiple comparisons. Five-fold cross-validation
gives fold-to-fold standard deviations between 0.018 and 0.032, which is four to six times
the 0.005 gap separating the top four models, and under cross-validation XGBoost narrowly
overtakes LightGBM. Our conclusion is that the leading models cannot be told apart on this
dataset, and that choosing a champion from a single split is not a reliable procedure.

A separate resampling study held the algorithm and its hyperparameters fixed and varied
only the sampling step. Class weighting was best (0.828), followed by random oversampling
(0.815), SMOTE (0.774), and random undersampling far behind (0.548). Undersampling
collapses because of how it works: balancing the classes by discarding majority rows
shrinks the training set from 170{,}235 rows to 568.

Further analysis shows that the tree-based models are well calibrated (expected calibration
error below 0.0025) while Logistic Regression is not (0.069), that the top decile of scored
transactions contains 92.6 percent of all frauds, which is a lift of 9.3 over random
review, and that under an amount-weighted cost model the tuned threshold is within 0.3
percent of the cost-optimal operating point. The final model is served through a FastAPI
backend and an interactive React web application with a 3D visualisation of the decision
surface.

\begin{keywords}
Credit Card Fraud Detection, Class Imbalance, PR-AUC, LightGBM, SMOTE, Resampling
Strategies, Bootstrap Confidence Intervals, McNemar's Test, Model Calibration,
Cost-Sensitive Thresholding
\end{keywords}
